\documentclass[11pt]{article}
\usepackage[utf8]{inputenc}
\usepackage[T1]{fontenc}
\usepackage{amsmath,amssymb,amsthm}
\usepackage[margin=1in]{geometry}
\usepackage{booktabs}
\usepackage[round]{natbib}
\setcitestyle{aysep={},yysep={;}}
\usepackage{hyperref}

\newtheorem{theorem}{Theorem}
\newtheorem{proposition}{Proposition}
\newtheorem{corollary}{Corollary}
\theoremstyle{remark}
\newtheorem{remark}{Remark}
\DeclareMathOperator{\Cov}{Cov}
\newcommand{\E}{\mathbb{E}}

\title{Advantage Alignment Does Not Preserve Mixed Nash Equilibria:\\
A Machine-Checked Counterexample and an Exact Covariance Characterization}
\author{Alejandro Panera Pascual (GalexY)\\ \small Universidad de Salamanca, Spain\\ \small \url{https://github.com/GalexY3}}
\date{November 2026}

\begin{document}
\maketitle

\begin{abstract}
Advantage Alignment (Duque et al., ICLR 2025) is a scalable opponent-shaping method whose Theorem~3 states that the alignment term vanishes at every Nash equilibrium, so that the update preserves Nash equilibria. We construct a two-player, two-stage, fully observed game whose unique Markov-perfect equilibrium is completely mixed. At that equilibrium the alignment field equals $(0,\beta\gamma^3/4,\beta\gamma^3/4)$ for one player. After one ascent step the other player has a strictly profitable unilateral deviation worth $\gamma(1+c)\tanh(\alpha\beta\gamma^3/8)>0$. We then prove that, at any state and for any pre-action weight $W$, the alignment term equals $W$ times the covariance, under the agent's own policy, between its score and the \emph{externality} of its action on the opponent's expected advantage. Nash indifference constrains each player's own-action marginal advantage, not this covariance. We derive the signed state weights that aggregate these conditional terms. Even with fixed local externalities, varying the information carried by the next state can cancel or reverse the update. Moreover, the update is not invariant to strategic equivalence: adding to one player's reward a term that depends only on the opponent's action leaves every unilateral deviation gain and the equilibrium unchanged, yet switches the equilibrium update on or off. Our counterexample is driven entirely by such a component. In normal-form games, the AA alignment term depends only on the main effect of the agent's action on the opponent's payoff (a welfare externality), whereas the first-order LOLA shaping term depends only on the interaction component (how the agent's action changes the opponent's incentives): the two methods shape through orthogonal components of the opponent's payoff. The same identity explains why the original argument does hold for pure-strategy profiles. We analyse two repairs. Marginalizing the opponent's current advantage removes the alignment term for every policy and reduces the method to naive policy gradient. Averaging the agent's past advantage over the opponent's action, as suggested by the original authors, restores preservation at completely mixed Markov-perfect equilibria but can remove the alignment signal on entire families of profiles. More generally, we characterize the equilibrium-preserving shaping terms in normal-form games: a local shaping term that is bilinear in the two players' payoffs vanishes at every interior Nash equilibrium if and only if it is \emph{incentive-gated}, reading the agent's own payoff only through its incentive gap or the opponent's payoff only through the opponent's incentive gap. AA and first-order LOLA are not gated; LookAhead and the second repair are. Hence no bilinear shaping term can both preserve interior equilibria and act at them: preservation forces the term to switch off exactly at equilibrium, leaving only its effect on local stability. The equilibrium witness, its update and profitable deviation, the conditional covariance identity and the failure of invariance under strategic equivalence have machine-checked proofs in Lean~4 with Mathlib. The state aggregation and the memory control are derived analytically and checked with exact arithmetic.
\end{abstract}

\section{Introduction}

Opponent shaping methods modify an agent's learning update to account for how other learning agents respond. LOLA \citep{foerster2018lola} differentiates through an opponent's learning step. Stable Opponent Shaping \citep{letcher2019sos} showed that LOLA can fail to preserve stable fixed points and proposed a correction. Later methods include POLA \citep{zhao2022pola}, COLA \citep{willi2022cola} and M-FOS \citep{lu2022mfos}. Advantage Alignment (AA) \citep{duque2025aa} replaces higher-order differentiation by a product of advantages that can be estimated with standard reinforcement learning, and reports strong results in social dilemmas.

Theorem~3 of \citet{duque2025aa} states that AA \emph{preserves Nash equilibria}: if the joint policy is a Nash equilibrium, ``the gradient contribution of the advantage alignment term is zero''. Its proof (Appendix A.8) relies on a lemma that advantages vanish at Nash for actions in the support. That lemma holds for each player's advantage \emph{averaged over the other player's action}. The alignment term in Equation~8, however, multiplies \emph{joint-action} advantages, which need not vanish at a mixed equilibrium.

\paragraph{Contributions.}
(i) An exact counterexample in which the AA update moves the unique, completely mixed Markov-perfect equilibrium of a two-stage game to a profile with a profitable deviation (Section~\ref{sec:counter}). (ii) An exact characterization of the alignment term as a covariance between the agent's score and the externality of its action on the opponent's advantage, with an if-and-only-if condition for vanishing at equilibrium and a corollary covering pure profiles (Section~\ref{sec:cov}). (iii) The aggregation of these conditional terms over histories and two controls showing that the equilibrium update depends on how past experience reaches the current state and on payoff components that are irrelevant to unilateral incentives, so AA is not invariant to strategic equivalence (Sections~\ref{sec:cov} and~\ref{sec:counter}). (iv) A decomposition showing that AA shapes the opponent through its welfare and first-order LOLA through its incentives (Section~\ref{sec:welfare}). (v) An analysis of two repairs (Section~\ref{sec:repair}). (vi) A characterization of all bilinear shaping terms that preserve interior Nash equilibria, which places AA, LOLA, LookAhead and both repairs in one classification and shows that equilibrium preservation and shaping at equilibrium are incompatible (Section~\ref{sec:class}). (vii) A Lean~4/Mathlib formalization of the main statements (Section~\ref{sec:lean}).

\paragraph{Context.} The first author of \citet{duque2025aa} publicly confirmed, in response to a question about this counterexample in the method's code repository, that Theorem~3 ``is too broad as stated'' and covers pure-strategy equilibria but not mixed equilibria in general, and suggested the second repair studied here \citep{duque2026issue}. The characterization, the analysis of both repairs and the machine-checked proofs are new to this note. We found no prior correction of Theorem~3 in the arXiv versions of \citet{duque2025aa}, in the titles of the works citing it indexed by OpenAlex, or in \citet{piche2026thesis}. The latter derives a related alignment term for a Stackelberg leader facing a soft best response, with past advantages marginalized over the agent's own action, and does not address Nash preservation.

\section{Related work}\label{sec:related}

\paragraph{Opponent shaping and learning awareness.} LOLA \citep{foerster2018lola} differentiates through the opponent's naive learning step. Its LookAhead variant drops the term that differentiates through the opponent's gradient, and Stable Opponent Shaping interpolates between the two so as to keep stable fixed points \citep{letcher2019sos}. Consistent and proximal refinements followed \citep{willi2022cola,zhao2022pola}, as did meta-learned shapers \citep{lu2022mfos,khan2024shaper}, Q-learning-aware shaping \citep{aghajohari2024loqa} and learning-aware policy gradients \citep{meulemans2025learningaware}. Meta-MAPG \citep{kim2021metamapg} adds a peer-learning term whose fixed-point displacement and basin selection are analysed by \citet{shcherbinin2026basin}. Advantage Alignment \citep{duque2025aa} shows that LOLA in REINFORCE form and LOQA are special cases of its update, so the scope of its Theorem~3 bears on that whole family.

\paragraph{Fixed points of learning dynamics.} In differentiable games, the Jacobian of the simultaneous gradient splits into potential and Hamiltonian parts, and the relation between stable fixed points and Nash equilibria depends on that split \citep{balduzzi2018mechanics,letcher2019mechanics}. Learning-aware dynamics can converge to points that are not Nash equilibria, as in Stackelberg learning \citep{fiez2020stackelberg}. \citet{letcher2019sos} state that a correction preserves fixed points when it vanishes where the simultaneous gradient vanishes. Section~\ref{sec:class} characterizes which corrections can do so, by the payoff components they read.

\paragraph{Decompositions and transformations of games.} We use the policy-weighted analysis-of-variance decomposition of payoffs and the notion of nonstrategic components of \citet{candogan2011}. \citet{tewolde2024transformations} characterize the game transformations that preserve Nash equilibria or best-response sets. Our strategic-equivalence control shows that the AA update is not invariant under such transformations.

\paragraph{Comparative statics of mixed equilibria and regulation.} In games without pure equilibria, changing one player's payoffs changes the opponent's equilibrium strategy rather than the player's own \citep{tsebelis1989}. \citet{weissing1991} discuss the conditions under which this holds. The remark in Section~\ref{sec:class} is the learning-dynamics counterpart of this observation. Finally, the good regulator theorem states that a successful regulator must model the system it regulates \citep{conant1970regulator}. AA regulates the opponent through its welfare, which Nash conditions do not constrain, and this is why it lacks the equilibrium guarantee.

\section{Setting}

Two players $i\in\{1,2\}$ act simultaneously in a Markov game with discount $\gamma\in(0,1)$. Player~1 takes $a_t$ and player~2 takes $b_t$ at state $s_t$. Policies $\pi_i(\cdot\mid s;\theta_i)$ are differentiable in their parameters and have full support (e.g.\ softmax). Write $S^1_t=\nabla_{\theta_1}\log\pi_1(a_t\mid s_t)$ and define exact advantages $A_i(s,a,b)=Q_i(s,a,b)-V_i(s)$, treated as constants when differentiating. Following Equation~8 of \citet{duque2025aa}, the AA update direction of player~1 is the ordinary policy gradient plus
\begin{equation}\label{eq:aa}
\beta\,\E\sum_{t\ge 0}\gamma^{t+1}\Big(\sum_{k<t}\gamma^{t-k}A_1(s_k,a_k,b_k)\Big)A_2(s_t,a_t,b_t)\,S^1_t .
\end{equation}
The bracket $W_t$ is measurable with respect to the history $h_t$ before the current actions.

\section{The alignment term is a covariance}\label{sec:cov}

\begin{proposition}[Covariance identity]\label{prop:cov}
Fix a history $h_t$ with current state $s=s_t$, and let $W=W(h_t)$. Define
$\mu_2(s,a)=\E_{b\sim\pi_2(\cdot\mid s)}A_2(s,a,b)$. Then
\[
\E\big[W A_2(s,a_t,b_t) S^1_t \,\big|\, h_t\big]
= W\,\Cov_{a\sim\pi_1(\cdot\mid s)}\big(\nabla_{\theta_1}\log\pi_1(a\mid s),\,\mu_2(s,a)\big).
\]
\end{proposition}
\begin{proof}
By independence of the simultaneous actions, the left side equals
$W\sum_a\pi_1(a\mid s)S^1(a)\mu_2(s,a)$. The score identity
$\sum_a\pi_1(a\mid s)S^1(a)=0$ lets us subtract any constant from $\mu_2$, in particular its mean.
\end{proof}

The function $a\mapsto\mu_2(s,a)-\E_a\mu_2(s,a)$ measures how player~1's current action changes player~2's expected advantage: an externality. The sign of $W$ determines whether the conditional term favors or discourages such actions.

\begin{corollary}[When the term vanishes]\label{cor:iff}
If $W\neq 0$, the conditional alignment term is zero if and only if the covariance in Proposition~\ref{prop:cov} is zero.
\end{corollary}

At a completely mixed equilibrium, Nash indifference states that $\E_{a\sim\pi_1}A_2(s,a,b)=0$ for every $b$ in the support of $\pi_2$. That is player~2's own-action marginal. It places no restriction on $\mu_2(s,\cdot)$, which averages over $b$ instead. For independent tabular softmax logits, the coordinate for action $a$ in the covariance is $\pi_1(a\mid s)[\mu_2(s,a)-\E_a\mu_2(s,a)]$. At a fixed history with $W\ne0$, constancy of $\mu_2$ on the support is thus necessary and sufficient for vanishing in this parameterization. With restricted parameterizations only its projection on the available score directions is observable.

\begin{proposition}[Aggregation over histories]\label{prop:aggregate}
In a finite stationary Markov game with bounded rewards, write $P$ for the on-policy state transition matrix, $u$ for the initial row distribution, and $C(s)$ for the score covariance in Proposition~\ref{prop:cov}. Define
\[
 M(s,j)=\E_{a,b\mid s}[A_1(s,a,b)\Pr(s_{t+1}=j\mid s,a,b)] .
\]
The full alignment term in Equation~\eqref{eq:aa} is $\beta\gamma\sum_s v(s)C(s)$, where
\[
 d=u(I-\gamma P)^{-1},\qquad v=\gamma^2 dM(I-\gamma^2P)^{-1}.
\]
\end{proposition}
\begin{proof}
Let $p_t(s)=\Pr(s_t=s)$ and $r_t(s)=\E[W_t\,\mathbf{1}_{s_t=s}]$. The recursion $W_{t+1}=\gamma(W_t+A_{1,t})$ gives $r_{t+1}=\gamma(r_tP+p_tM)$ with $r_0=0$. Summing the bounded discounted series gives $d=\sum_t\gamma^tp_t$ and $v=\sum_t\gamma^tr_t=\gamma^2(vP+dM)$. Applying Proposition~\ref{prop:cov} at each history yields the claimed expression.
\end{proof}
For disjoint tabular logit blocks and $\beta>0$, the full field of alignment vanishes exactly when each state has either zero signed weight $v(s)$ or constant $\mu_2(s,\cdot)$ on its support. Shared parameters permit further cancellations between states. If transitions are action-independent on the policy's Markov state space, $M=0$ by centering of the exact advantages. This also explains why an externality at a single stage alone cannot establish a nonzero update.

\begin{corollary}[Pure profiles]\label{cor:pure}
For a score defined on the sampled support, if $\pi_1(\cdot\mid s)$ is a point mass, the covariance is zero for every opponent advantage. Hence the alignment term vanishes at pure profiles, consistently with the scope stated by \citet{duque2026issue}. This is a boundary statement; finite softmax logits realize only interior policies.
\end{corollary}

\section{A counterexample}\label{sec:counter}

Both players choose signed actions in $\{-1,+1\}$ in two stages: $x,z$ at the start and $f,g$ afterwards. The public next state is $z$, and the game then ends. For $c>0$ the rewards are:
\begin{center}
\begin{tabular}{lcc}\toprule
Stage & Player 1 & Player 2\\\midrule
0 & $z(1+cx)$ & $-cxz$\\
1 (state $z$) & $cfg$ & $f(z-cg)$\\\bottomrule
\end{tabular}
\end{center}
Each player has three Bernoulli logits (start, state $+1$, state $-1$). Let $\pi^*$ be the uniform profile (all logits zero).

\begin{proposition}\label{prop:nash}
$\pi^*$ is a Markov-perfect equilibrium. All state values are zero, exact joint advantages equal the stage rewards, and no deterministic history-dependent unilateral deviation yields a positive return. It is the unique Markov-perfect equilibrium.
\end{proposition}

\begin{proof}
At either terminal state let $m_f=\E f$ and $m_g=\E g$. The payoff gains from choosing $+1$ instead of $-1$ are $2c m_g$ for player~1 and $-2c m_f$ for player~2. If $m_g>0$, player~1 strictly chooses $+1$ and player~2 then strictly chooses $-1$, a contradiction. The case $m_g<0$ reverses the signs. Thus $m_g=0$, and mixing by player~2 requires $m_f=0$. The terminal values are consequently zero. At the start the gains have the identical form $2c\,\E z$ and $-2c\,\E x$, forcing both means to zero again. Backward induction gives the unique Markov-perfect equilibrium. Against the uniform opponent, either player's conditional expected terminal reward is zero for each own action, even given the past. At the first stage the same statement holds after including continuation value. Hence even a history-dependent deviation cannot improve the return.
\end{proof}
This profile also satisfies the Boltzmann action-value relation used in deriving AA: each player's marginal action values coincide state by state.

\begin{proposition}\label{prop:field}
At $\pi^*$ the AA direction of Equation~\eqref{eq:aa}, including the ordinary policy gradient, is $F_1=(0,\beta\gamma^3/4,\beta\gamma^3/4)$ and $F_2=(0,0,0)$.
\end{proposition}
\begin{proof}[Sketch]
Ordinary gradients vanish by indifference. For player~1 at $t=1$ the integrand is $\beta\gamma^3\,z(1+cx)\,f(z-cg)\,f/2=\beta\gamma^3(1+cx)(1-czg)/2$, whose expectation given $z$ is $\beta\gamma^3/2$. Each terminal state has probability $1/2$. For player~2 the analogous product has mean zero. The Lean proof sums all 16 trajectories.
\end{proof}

\begin{proposition}\label{prop:dev}
Let $\alpha,\beta,\gamma,c>0$ and apply the simultaneous ascent step $\theta'=\theta^*+\alpha F$. Player~2's on-policy return remains $0$. The deviation ``$z=+1$, then $g=-1$'' earns
\[
\gamma(1+c)\big(2\sigma(\alpha\beta\gamma^3/4)-1\big)=\gamma(1+c)\tanh(\alpha\beta\gamma^3/8)>0 .
\]
Hence $\theta'$ is not a Nash equilibrium.
\end{proposition}

For $\gamma=0.9$, $\beta=0.3$, $\alpha=0.1$ and $c=0.1$, each nonzero field component is $0.054675$ and the deviation gain is $\approx 0.0027064$. The update does not move within a set of equilibria: it creates exploitability. In the notation of Proposition~\ref{prop:cov}, player~1's terminal action shifts player~2's advantage by $f\,z$, which covaries with player~1's score.

\subsection{Two controls that separate dynamics from incentives}
First replace the transition $s=z$ by $\Pr(s\mid z)=(1+\eta s z)/2$ with $\eta\in[-1,1]$, retaining the terminal payoff $f(s-cg)$. Terminal best-response differences and continuation values at the uniform profile are unchanged, so the preceding backward induction still applies. At either terminal state $C(s)=s/2$ in its logit coordinate, whereas $\E[A_{1,0}\mathbf{1}_{s_1=s}]=\eta s/2$. The two terminal field coordinates are now $\beta\gamma^3\eta/4$. Forgetting the initial choice ($\eta=0$) cancels the update; reversing its encoding while keeping the reward table fixed ($\eta=-1$) reverses its direction. This intervention changes the physical link between the initial choice and the subsequent payoff.

Second, return to $s=z$ and replace player~1's initial payment by $cxz+\lambda z$ for $\lambda\in\mathbb{R}$. At a fixed opponent strategy, the change from $\lambda=0$ adds $\lambda\E z$ to the player's return for \emph{every} own strategy. All unilateral deviation gains and equilibrium sets therefore coincide across the family. This is a nonstrategic component in the strategic-form sense of \citet{candogan2011}. Nevertheless, the two terminal AA coordinates at the uniform equilibrium become $\beta\gamma^3\lambda/4$, by the same trajectory calculation as in Proposition~\ref{prop:field}. By contrast, merely delaying $\lambda z$ to a terminal payment of $\lambda z/\gamma$ preserves discounted trajectory returns and all exact joint advantages: at the terminal state it adds a state-dependent constant to both $Q_1$ and $V_1$. These controls separate sensitivity to a payoff component that is irrelevant to unilateral deviations from sensitivity to the accounting date of that component. Strategic equivalence need not preserve welfare or the value of changing an opponent's future policy.

\section{Welfare shaping versus incentive shaping}\label{sec:welfare}

Consider a single state with finite action sets, the agent's softmax policy $p$ and the opponent's policy $q$. Decompose the opponent's payoff in the policy-weighted two-way analysis-of-variance form
\[
u_2(a,b)=\kappa+n_2(a)+s_2(b)+i_2(a,b),
\]
with $\E_p n_2=0$, $\E_q s_2=0$, $\E_{b\sim q}i_2(a,b)=0$ for every $a$ and $\E_{a\sim p}i_2(a,b)=0$ for every $b$. The term $n_2$ is the welfare externality of the agent's action: it does not affect the opponent's incentives and is a nonstrategic component for the opponent \citep{candogan2011}. The term $i_2$ is the interaction: it determines how the agent's action changes the opponent's relative payoffs across its own actions. Write $J_p=\mathrm{diag}(p)-pp^\top$ for the softmax Jacobian.

\begin{proposition}[AA reads welfare, LOLA reads incentives]\label{prop:welfare}
(i) The AA alignment term at the state, $W\,\E_{a,b}[(u_2(a,b)-V_2)\nabla_{\theta_1}\log p(a)]$, equals $W J_p\mu_2$ with $\mu_2(a)=\E_{b\sim q}u_2(a,b)$. It is unchanged by adding any $s_2$ or any $q$-centered $i_2$, and adding $n_2$ changes it by $W J_p n_2$.
(ii) The first-order LOLA shaping term of the agent \citep{foerster2018lola}, $(\nabla_{\theta_1}\nabla_{\theta_2}V_2)^\top\nabla_{\theta_2}V_1=J_p u_2 J_q g$ with $g=\nabla_{\theta_2}V_1$, is unchanged by adding any $n_2$ or $s_2$, and adding $i_2$ changes it by $J_p i_2 J_q g$.
\end{proposition}
\begin{proof}
(i) Proposition~\ref{prop:cov} gives $W J_p\mu_2$; adding $s_2$ shifts $\mu_2$ by a constant, which $J_p$ annihilates because $J_p\mathbf 1=0$; a $q$-centered $i_2$ leaves $\mu_2$ unchanged. (ii) $J_q\mathbf 1=0$ removes every term constant in $b$, including $n_2(a)$, and $\mathbf 1^\top J_p=0$ removes every term constant in $a$, including $s_2(b)$.
\end{proof}

The two components are separable. In a ``gift'' game where the agent's action only transfers welfare ($u_2(a,b)=\mathbf 1[a=1]$), the AA term at the uniform profile is $1/4$ while the LOLA term is $0$ for every $g$. In an ``incentive'' game with a pure interaction ($u_2(a,b)=ab$ for signed actions), the AA term is $0$ while the LOLA term is $1/4$ for $g=e_1$. AA therefore shapes like conditional reciprocity: the past weight records how the opponent's actions affected the agent, and the current term favors actions that benefit the opponent. First-order LOLA shapes like incentive design: it favors actions that change what is best for the opponent, as sanctions and rewards do. Neither reads the opponent's own main effect $s_2$. This also explains the strategic-equivalence control of Section~\ref{sec:counter}: Nash conditions are blind to welfare externalities, and AA is built from them.

\section{Two repairs}\label{sec:repair}

\begin{proposition}[Marginalizing the current opponent advantage trivializes AA]\label{prop:marg}
Replacing $A_2(s,a_t,b_t)$ by $\bar A_2(s,b_t)=\E_{a\sim\pi_1}A_2(s,a,b_t)$ makes the alignment term zero for \emph{every} policy profile.
\end{proposition}
\begin{proof}
$\bar A_2$ does not depend on $a_t$, so the conditional expectation factorizes and contains $\sum_a\pi_1(a\mid s)S^1(a)=0$.
\end{proof}

This removes the failure only by removing the method: the update becomes naive policy gradient.

\begin{proposition}[Marginalizing the past weight]\label{prop:past}
Replace each $A_1(s_k,a_k,b_k)$ in $W_t$ by $\tilde A_1(s_k,a_k)=\E_{b\sim\pi_2}A_1(s_k,a_k,b)$, as suggested by \citet{duque2026issue}. At a Markov-perfect equilibrium where $\pi_1$ is completely mixed at every state, $\tilde A_1\equiv 0$ by indifference, so the repaired alignment term vanishes, and the profile is a fixed point of the repaired update.
\end{proposition}

\begin{remark}
The repaired weight depends only on the quality of the agent's own past actions. If those actions do not influence future states, its expectation is zero and independent of the current covariance. In the game of Section~\ref{sec:counter}, player~1's repaired alignment term is therefore identically zero at \emph{every} profile, while player~2's is not. We confirmed this numerically on 2000 random profiles (largest absolute value $2.5\times10^{-17}$ for player~1). Preserving mixed equilibria in this way can thus suppress the alignment signal, including its dependence on the opponent's past action that motivates reciprocity. Section~\ref{sec:class} shows that this is not an accident of the repair: within bilinear terms, preserving every interior equilibrium requires the term to vanish there. \citet{piche2026thesis} marginalizes in the opposite direction, to isolate the opponent's contribution, for a different (Stackelberg) objective.
\end{remark}

\section{Which shaping terms preserve equilibria}\label{sec:class}

Keep the single-state setting of Section~\ref{sec:welfare} and also decompose the agent's own payoff, $u_1=\kappa_1+r_1(a)+t_1(b)+j_1(a,b)$, with the same centering. At interior policies $(p,q)$, the profile is a Nash equilibrium of $(u_1,u_2)$ if and only if $r_1=0$ and $s_2=0$: each player is indifferent among its actions against the other's policy. Call $r_1$ and $s_2$ the \emph{incentive gaps}. A local shaping term of the agent is a map $T(u_1,u_2)\in\mathbb R^{A}$, bilinear in the two payoff functions for fixed $(p,q)$, added to the policy gradient $J_p r_1$ with weight $\beta$.

\begin{theorem}[Equilibrium-preserving shaping is incentive-gated]\label{thm:gated}
$T$ vanishes at every game for which $(p,q)$ is an interior Nash equilibrium if and only if there are bilinear maps $B_1,B_2$ with
\[
T(u_1,u_2)=B_1(r_1,u_2)+B_2(u_1,s_2)\quad\text{for all }u_1,u_2 .
\]
Equivalently, for $\beta\neq0$ the shaped update $J_pr_1+\beta T$ has every interior Nash equilibrium as a fixed point if and only if $T$ is gated.
\end{theorem}
\begin{proof}
Sufficiency is immediate. For necessity, choose linear generalized inverses $L_1,L_2$ with $r_1L_1r_1=r_1$ and $s_2L_2s_2=s_2$, and write $u_1=g_1+L_1r_1(u_1)$ and $u_2=g_2+L_2s_2(u_2)$ with $r_1(g_1)=0$ and $s_2(g_2)=0$. Bilinearity gives $T(u_1,u_2)=T(g_1,g_2)+T(L_1r_1u_1,u_2)+T(g_1,L_2s_2u_2)$, and the first term is zero by hypothesis. The second term is $B_1(r_1,u_2)$, and the third is $B_2(u_1,s_2)$ because $g_1=u_1-L_1r_1u_1$ is linear in $u_1$.
\end{proof}

\begin{center}
\begin{tabular}{lll}\toprule
Term & Reads & Gated?\\\midrule
AA, Eq.~\eqref{eq:aa} & joint past advantage $\times$ welfare $n_2$ & no\\
First-order LOLA & interaction $i_2$ and $t_1$ & no\\
LookAhead (LOLA without $\chi$; \citealp{letcher2019sos}) & opponent gap $s_2$ & yes\\
AA, current advantage marginalized (Prop.~\ref{prop:marg}) & nothing (identically zero) & yes\\
AA, past weight marginalized (Prop.~\ref{prop:past}) & own gap $r_1$ & yes\\\bottomrule
\end{tabular}
\end{center}

The failure of LOLA to preserve fixed points was established by \citet{letcher2019sos}, and \citet{shcherbinin2026basin} show an $O(\lambda)$ displacement for the peer-learning term of Meta-MAPG. Theorem~\ref{thm:gated} places these and the AA counterexample in one classification by payoff components.

\begin{corollary}[No free shaping at equilibrium]\label{cor:nofree}
If a bilinear shaping term preserves every interior Nash equilibrium, it is zero at all of them. At an equilibrium it can change only the local Jacobian of the update, hence stability and selection among equilibria, not their location. Every term that acts at equilibrium by reading welfare ($n_2$), interaction ($i_2$) or the opponent's effect on the agent ($t_1$) displaces some interior equilibrium.
\end{corollary}

This sharpens the question left by the repairs: equilibrium-preserving variants of AA can keep their effect away from equilibrium and their effect on stability, but not their effect at the equilibrium itself.

\begin{remark}[Where the displacement goes]
Suppose only the agent shapes and the opponent follows its ordinary policy gradient. The opponent's fixed-point condition $J_q(u_2^\top p)=0$ at an interior $q$ says only that $u_2^\top p$ is constant. It involves neither $q$, $\beta$ nor $T$. Whenever it determines $p$ uniquely, which is generic when the agent has no more actions than the opponent, every fixed point of the shaped dynamics has the same $p$: shaping moves only the opponent's equilibrium strategy. To first order in a square game this displacement is $-\beta H_{12}^{-1}T$, where $H_{12}=J_pu_1J_q$ in reduced logits. This is the dynamic counterpart of the classic observation that in games without pure equilibria a change in one player's payoffs changes the opponent's equilibrium strategy, not the player's own \citep{tsebelis1989}; \citet{weissing1991} discuss when this holds. In numerical checks with 500 random games per size, the agent's displacement was zero to machine precision in square games ($2\times2$ to $4\times4$) and in $2\times3$ and $3\times4$ games. When the agent has more actions than the opponent ($3\times2$, $4\times3$), the shaped dynamics usually had no nearby interior fixed point.
\end{remark}

\section{Machine verification}\label{sec:lean}

We formalized the statements in Lean~4 (v4.36.0) \citep{demoura2021lean4} with Mathlib \citep{mathlib2020}. The development defines the game, Bernoulli-logit policies with their analytic score derivatives, exact values and joint advantages, and Equation~\eqref{eq:aa} as an exact expectation over all trajectories. It proves:
\begin{itemize}
\item Proposition~\ref{prop:cov} in finite-action form, Corollary~\ref{cor:iff}, Corollary~\ref{cor:pure}, and the score identity (softmax, and Bernoulli with derivatives);
\item the equilibrium part of Proposition~\ref{prop:nash} for all deterministic history-dependent deviations of either player, Proposition~\ref{prop:field}, and the fact that the step in Proposition~\ref{prop:dev} is exactly $\alpha F$;
\item Proposition~\ref{prop:dev} for all $\alpha,\beta,\gamma,c>0$;
\item Proposition~\ref{prop:marg}, and the vanishing part of Proposition~\ref{prop:past} given indifference;
\item Proposition~\ref{prop:welfare} (both parts, given the softmax Jacobian) and the gift and incentive examples;
\item Theorem~\ref{thm:gated} and Corollary~\ref{cor:nofree} for bilinear maps between arbitrary vector spaces over any field, with arbitrary linear incentive-gap maps (the game is an instance); and the exact-anchoring statement of the remark in Section~\ref{sec:class}, with uniqueness of the indifference solution as a hypothesis;
\item the strategic-equivalence control for every $\lambda$: the $\lambda z$ term shifts player~1's return by $\lambda\,\E z$ for every deterministic history-dependent policy against every opponent profile, so deviation gains do not depend on $\lambda$; the field at the uniform equilibrium is $(0,\beta\gamma^3\lambda/4,\beta\gamma^3\lambda/4)$; $\lambda=1$ recovers the original operator; and for $\beta,\gamma>0$ the update vanishes at $\lambda=0$ but not at any $\lambda\neq0$.
\end{itemize}
\paragraph{Confirmation with the original implementation.} We also executed the authors' unmodified advantage and loss functions (\texttt{compute\_gae\_advantages}, \texttt{aa\_terms}, \texttt{aa\_loss}, \texttt{actor\_losses}; code revision 2065e43), loaded with checksum verification, on the $\lambda$ family with an exact critic. In all 80 configurations with exact one-step advantages (two loss modes, proximal or not, normalized or not, single and double precision), the update vanishes at $\lambda=0$ and leaves the profile unexploitable, whereas for every $\lambda\neq0$ it moves the terminal logits with the sign of $\lambda$ and creates a profitable deviation; without normalization the update is proportional to $\lambda$. With generalized advantage estimation ($\lambda_{\mathrm{GAE}}>0$) future rewards enter the weight and produce an additional small movement even at $\lambda=0$.

\texttt{\#print axioms} reports only \texttt{propext}, \texttt{Classical.choice} and \texttt{Quot.sound} for every main theorem, and the development contains no \texttt{sorry}. Not formalized: the derivative identity $\nabla\,\mathrm{softmax}=J_p$ (standard), the uniqueness argument, Proposition~\ref{prop:aggregate}, the memory control ($\eta$) and the delayed-payment variant, mixed deviations (convex combinations of deterministic ones), the numerical remark of Section~\ref{sec:repair}, the classification table of Section~\ref{sec:class} for LookAhead and the two repairs (checked numerically on 2000 random games per size, $|T|\le 2\times10^{-16}$), the genericity of the uniqueness hypothesis, the first-order displacement formula, and sampled generalized advantage estimators. The current development has 37 audited theorems.

\section{Discussion and limitations}

Our results concern the exact-advantage operator of Equation~8. They do not show that AA performs poorly. Moving away from a Nash equilibrium can be desirable in a social dilemma, and the alignment force we characterize is precisely what AA intends. What fails is the claimed equilibrium guarantee. Implementations with learned critics or sampled advantage estimators define different random operators and are not covered by the formal statements. A reading of Appendix~A.8, Equation~66, in which the inner expectation is an unconditional scalar, would define a different operator.

\paragraph{Reproducibility.} The Lean project, executable checks of the equilibrium witness, the state aggregation and the strategic-equivalence control, and the preregistration of the formal statements are available at \url{https://github.com/GalexY3/aa-nash-counterexample}.

\paragraph{Acknowledgements and contributions.} The research idea is the author's, who raised the counterexample with the authors of \citet{duque2025aa} in \citet{duque2026issue}. Advantage Alignment, its Equation~8 and Theorem~3 are entirely the work of \citet{duque2025aa}; this note only examines the scope of that theorem. We thank Juan Agustin Duque for his public response, which confirmed the scope issue and proposed the repair analysed in Proposition~\ref{prop:past}. Alejandro Panera Pascual conceived and directed the investigation; AI tools were used under his direction for derivations, computational checks, the Lean formalization and manuscript preparation.

\begin{thebibliography}{99}
\setlength{\itemsep}{2pt}

\bibitem[Aghajohari et~al.(2024)]{aghajohari2024loqa} Aghajohari, Milad, Juan Agustin Duque, Tim Cooijmans, and Aaron Courville. 2024. ``LOQA: Learning with Opponent Q-Learning Awareness.'' In \emph{International Conference on Learning Representations (ICLR)}. arXiv:2405.01035.

\bibitem[Balduzzi et~al.(2018)]{balduzzi2018mechanics} Balduzzi, David, S\'ebastien Racani\`ere, James Martens, Jakob Foerster, Karl Tuyls, and Thore Graepel. 2018. ``The Mechanics of n-Player Differentiable Games.'' In \emph{Proceedings of the 35th International Conference on Machine Learning}, Proceedings of Machine Learning Research 80. \url{https://proceedings.mlr.press/v80/balduzzi18a.html}.

\bibitem[Candogan et~al.(2011)]{candogan2011} Candogan, Ozan, Ishai Menache, Asuman Ozdaglar, and Pablo A. Parrilo. 2011. ``Flows and Decompositions of Games: Harmonic and Potential Games.'' \emph{Mathematics of Operations Research} 36(3):474--503. doi:10.1287/moor.1110.0500.

\bibitem[Conant and Ashby(1970)]{conant1970regulator} Conant, Roger C., and W. Ross Ashby. 1970. ``Every Good Regulator of a System Must Be a Model of That System.'' \emph{International Journal of Systems Science} 1(2):89--97. doi:10.1080/00207727008920220.

\bibitem[de~Moura and Ullrich(2021)]{demoura2021lean4} de~Moura, Leonardo, and Sebastian Ullrich. 2021. ``The Lean 4 Theorem Prover and Programming Language.'' Pp. 625--35 in \emph{Automated Deduction -- CADE 28}, Lecture Notes in Computer Science 12699. Cham: Springer. doi:10.1007/978-3-030-79876-5\_37.

\bibitem[Duque(2026)]{duque2026issue} Duque, Juan Agustin. 2026. ``Reply to `Question about Theorem 3: Joint Advantages at a Mixed Nash Equilibrium'.'' GitHub issue \#12, jduquevan/advantage-alignment, November 16. \url{https://github.com/jduquevan/advantage-alignment/issues/12}.

\bibitem[Duque et~al.(2025)]{duque2025aa} Duque, Juan Agustin, Milad Aghajohari, Tim Cooijmans, Razvan Ciuca, Tianyu Zhang, Gauthier Gidel, and Aaron Courville. 2025. ``Advantage Alignment Algorithms.'' In \emph{International Conference on Learning Representations (ICLR)}. arXiv:2406.14662.

\bibitem[Fiez et~al.(2020)]{fiez2020stackelberg} Fiez, Tanner, Benjamin Chasnov, and Lillian Ratliff. 2020. ``Implicit Learning Dynamics in Stackelberg Games: Equilibria Characterization, Convergence Analysis, and Empirical Study.'' In \emph{Proceedings of the 37th International Conference on Machine Learning}, Proceedings of Machine Learning Research 119. \url{https://proceedings.mlr.press/v119/fiez20a.html}.

\bibitem[Foerster et~al.(2018)]{foerster2018lola} Foerster, Jakob, Richard Y. Chen, Maruan Al-Shedivat, Shimon Whiteson, Pieter Abbeel, and Igor Mordatch. 2018. ``Learning with Opponent-Learning Awareness.'' In \emph{Proceedings of the 17th International Conference on Autonomous Agents and MultiAgent Systems (AAMAS)}. arXiv:1709.04326.

\bibitem[Khan et~al.(2024)]{khan2024shaper} Khan, Akbir, Timon Willi, Newton Kwan, Andrea Tacchetti, Chris Lu, Edward Grefenstette, Tim Rockt\"aschel, and Jakob Foerster. 2024. ``Scaling Opponent Shaping to High Dimensional Games.'' In \emph{Proceedings of the 23rd International Conference on Autonomous Agents and MultiAgent Systems (AAMAS)}. arXiv:2312.12568.

\bibitem[Kim et~al.(2021)]{kim2021metamapg} Kim, Dong-Ki, Miao Liu, Matthew Riemer, Chuangchuang Sun, Marwa Abdulhai, Golnaz Habibi, Sebastian Lopez-Cot, Gerald Tesauro, and Jonathan P. How. 2021. ``A Policy Gradient Algorithm for Learning to Learn in Multiagent Reinforcement Learning.'' In \emph{Proceedings of the 38th International Conference on Machine Learning}, Proceedings of Machine Learning Research 139. arXiv:2011.00382.

\bibitem[Letcher et~al.(2019a)]{letcher2019mechanics} Letcher, Alistair, David Balduzzi, S\'ebastien Racani\`ere, James Martens, Jakob Foerster, Karl Tuyls, and Thore Graepel. 2019a. ``Differentiable Game Mechanics.'' \emph{Journal of Machine Learning Research} 20(84):1--40. arXiv:1905.04926.

\bibitem[Letcher et~al.(2019b)]{letcher2019sos} Letcher, Alistair, Jakob Foerster, David Balduzzi, Tim Rockt\"aschel, and Shimon Whiteson. 2019b. ``Stable Opponent Shaping in Differentiable Games.'' In \emph{International Conference on Learning Representations (ICLR)}. arXiv:1811.08469.

\bibitem[Lu et~al.(2022)]{lu2022mfos} Lu, Chris, Timon Willi, Christian Schroeder de~Witt, and Jakob Foerster. 2022. ``Model-Free Opponent Shaping.'' In \emph{Proceedings of the 39th International Conference on Machine Learning}, Proceedings of Machine Learning Research 162. arXiv:2205.01447.

\bibitem[mathlib Community(2020)]{mathlib2020} The mathlib Community. 2020. ``The Lean Mathematical Library.'' Pp. 367--81 in \emph{Proceedings of the 9th ACM SIGPLAN International Conference on Certified Programs and Proofs (CPP 2020)}. New York: ACM. doi:10.1145/3372885.3373824.

\bibitem[Meulemans et~al.(2025)]{meulemans2025learningaware} Meulemans, Alexander, Seijin Kobayashi, Johannes von Oswald, Nino Scherrer, Eric Elmoznino, Blake Richards, Guillaume Lajoie, Blaise Ag\"uera y Arcas, and Jo\~ao Sacramento. 2025. ``Multi-Agent Cooperation through Learning-Aware Policy Gradients.'' In \emph{International Conference on Learning Representations (ICLR)}. arXiv:2410.18636.

\bibitem[Pich\'e(2026)]{piche2026thesis} Pich\'e, D. 2026. ``Advantage Alignment and Large Language Models.'' Master's thesis, Universit\'e de Montr\'eal. \url{https://hdl.handle.net/1866/44392}.

\bibitem[Shcherbinin et~al.(2026)]{shcherbinin2026basin} Shcherbinin, Yevhen, Arina Redina, Maxim Kalpin, and Vlad Kochetov. 2026. ``Equilibrium Selection in Multi-Agent Policy Gradients via Opponent-Aware Basin Entry.'' arXiv:2605.18078.

\bibitem[Tewolde and Conitzer(2024)]{tewolde2024transformations} Tewolde, Emanuel, and Vincent Conitzer. 2024. ``Game Transformations That Preserve Nash Equilibria or Best-Response Sets.'' In \emph{Proceedings of the 33rd International Joint Conference on Artificial Intelligence (IJCAI)}. arXiv:2111.00076.

\bibitem[Tsebelis(1989)]{tsebelis1989} Tsebelis, George. 1989. ``The Abuse of Probability in Political Analysis: The Robinson Crusoe Fallacy.'' \emph{American Political Science Review} 83(1):77--91.

\bibitem[Weissing and Ostrom(1991)]{weissing1991} Weissing, Franz, and Elinor Ostrom. 1991. ``Crime and Punishment: Further Reflections on the Counterintuitive Results of Mixed Equilibria Games.'' \emph{Journal of Theoretical Politics} 3(3):343--50.

\bibitem[Willi et~al.(2022)]{willi2022cola} Willi, Timon, Alistair Letcher, Johannes Treutlein, and Jakob Foerster. 2022. ``COLA: Consistent Learning with Opponent-Learning Awareness.'' In \emph{Proceedings of the 39th International Conference on Machine Learning}, Proceedings of Machine Learning Research 162. arXiv:2203.04098.

\bibitem[Zhao et~al.(2022)]{zhao2022pola} Zhao, Stephen, Chris Lu, Roger B. Grosse, and Jakob Foerster. 2022. ``Proximal Learning with Opponent-Learning Awareness.'' In \emph{Advances in Neural Information Processing Systems 35 (NeurIPS)}. arXiv:2210.10125.
\end{thebibliography}
\end{document}
